\documentclass[11pt,a4paper]{article}
\usepackage{amsmath,amssymb}
\usepackage{geometry}
\usepackage{hyperref}
\usepackage{listings}
\usepackage{xcolor}

\geometry{margin=1in}

\title{\textbf{Sudoku Solver Pseudocode}\\
\large Minigrid Permutations, Compatibility Graphs, and Exact Support Search}
\author{}
\date{}

\lstdefinelanguage{Pseudo}{
    morekeywords={procedure,function,input,output,for,each,in,if,then,else,while,do,return,break,continue,let,match,case,repeat,until,true,false,None,Some},
    sensitive=false,
    morecomment=[l]{//}
}

\lstset{
    language=Pseudo,
    basicstyle=\ttfamily\small,
    keywordstyle=\color{blue}\bfseries,
    commentstyle=\color{gray},
    breaklines=true,
    frame=single,
    columns=fullflexible,
    keepspaces=true,
    showstringspaces=false
}

\begin{document}

\maketitle

\section{Notation}

The solver works with an $N \times N$ Sudoku board and $K \times K$ minigrids. For the standard puzzle, $N = 9$ and $K = 3$. The implementation treats each minigrid completion as a graph vertex. A full solution is an assignment of one compatible permutation to each minigrid.

\begin{itemize}
    \item $B$ is the input board. A zero cell is empty.
    \item $M$ stores row, column, box, and per-cell conflict masks.
    \item $P_i$ is the set of generated permutations for minigrid $i$.
    \item $G$ is the compatibility graph over all minigrid permutations.
    \item $D_i$ is the current search domain of permutation ids for minigrid $i$.
    \item $A_i$ is the selected permutation id for minigrid $i$, if assigned.
\end{itemize}

\section{Complete Solver Pipeline}

This algorithm corresponds to \texttt{SudokuSolver::solve\_with\_stats} in \texttt{src/solver/mod.rs}. It is the top-level pipeline used by the full reporting path.

\begin{lstlisting}[caption={Complete solver pipeline}]
procedure SolveWithStats(board B)
    start total timer

    // Phase 1: fixed conflicts from the input board
    M <- GenerateMasks(B)

    // Phase 1.5: deterministic preprocessing
    H <- copy(B)
    filled <- ApplyHiddenSingles(H, M)

    // Phase 2: local minigrid search
    for each minigrid id i in 0 .. N - 1 do in parallel
        P_i <- GenerateMinigridPermutations(H, M, i)
    end for

    // Phase 3: graph of compatible local completions
    G <- NewGraph(P_0, P_1, ..., P_{N-1})
    BuildCompatibilityEdges(G)

    // Phase 4: exact support pruning
    (removed, configurations) <- ExactSupportPrune(G)

    // Phase 5: reconstruct boards from full configurations
    solutions <- ExtractSolutions(G, configurations)
    class <- Classify(solutions)

    return SolveReport(solutions, statistics, class)
end procedure
\end{lstlisting}

\section{Mask Initialization}

This algorithm corresponds to \texttt{Masks::generate} in \texttt{src/types/masks.rs}. The masks reject invalid givens and provide constant-time conflict checks for later stages.

\begin{lstlisting}[caption={Generate row, column, box, and conflict masks}]
procedure GenerateMasks(board B)
    initialize rows[0 .. N-1] as empty masks
    initialize cols[0 .. N-1] as empty masks
    initialize boxes[0 .. N-1] as empty masks
    initialize conflict[0 .. N-1][0 .. N-1] as empty masks

    for r in 0 .. N - 1 do
        for c in 0 .. N - 1 do
            d <- B[r][c]
            if d = 0 then
                continue
            end if

            b <- BoxIndex(r, c)
            if rows[r] contains d or cols[c] contains d or boxes[b] contains d then
                error "invalid board: duplicate given"
            end if

            add d to rows[r]
            add d to cols[c]
            add d to boxes[b]
        end for
    end for

    for r in 0 .. N - 1 do
        for c in 0 .. N - 1 do
            b <- BoxIndex(r, c)
            conflict[r][c] <- rows[r] union cols[c] union boxes[b]
        end for
    end for

    return Masks(rows, cols, boxes, conflict)
end procedure
\end{lstlisting}

\section{Hidden Single Deduction}

This algorithm corresponds to \texttt{apply\_hidden\_singles} in \texttt{src/solver/heuristics.rs}. It fills digits that have one legal position in a row, column, or minigrid.

\begin{lstlisting}[caption={Hidden single deduction}]
procedure ApplyHiddenSingles(board B, masks M)
    total_filled <- 0

    repeat
        changed <- false

        for d in 1 .. N do
            for each row r do
                if FillIfOnlyLegalCell(B, M, d, cells in row r) then
                    total_filled <- total_filled + 1
                    changed <- true
                end if
            end for

            for each column c do
                if FillIfOnlyLegalCell(B, M, d, cells in column c) then
                    total_filled <- total_filled + 1
                    changed <- true
                end if
            end for

            for each box b do
                if FillIfOnlyLegalCell(B, M, d, cells in box b) then
                    total_filled <- total_filled + 1
                    changed <- true
                end if
            end for
        end for
    until changed = false

    return total_filled
end procedure

function FillIfOnlyLegalCell(board B, masks M, digit d, cell_set S)
    count <- 0
    possible <- None

    for each cell (r, c) in S do
        if B[r][c] = 0 and d not in M.conflict[r][c] then
            possible <- Some(r, c)
            count <- count + 1
        end if
    end for

    if count = 1 then
        (r, c) <- possible
        B[r][c] <- d
        UpdateMasksAfterPlacement(M, r, c, d)
        return true
    end if

    return false
end function
\end{lstlisting}

\begin{lstlisting}[caption={Update masks after a deterministic placement}]
procedure UpdateMasksAfterPlacement(masks M, row r, column c, digit d)
    b <- BoxIndex(r, c)
    add d to M.rows[r]
    add d to M.cols[c]
    add d to M.boxes[b]

    for each column x in 0 .. N - 1 do
        bx <- BoxIndex(r, x)
        M.conflict[r][x] <- M.rows[r] union M.cols[x] union M.boxes[bx]
    end for

    for each row y in 0 .. N - 1 do
        by <- BoxIndex(y, c)
        M.conflict[y][c] <- M.rows[y] union M.cols[c] union M.boxes[by]
    end for

    for each cell (y, x) inside box b do
        bx <- BoxIndex(y, x)
        M.conflict[y][x] <- M.rows[y] union M.cols[x] union M.boxes[bx]
    end for
end procedure
\end{lstlisting}

\section{Minigrid Permutation Generation}

This algorithm corresponds to \texttt{generate\_all\_permutations} and \texttt{PermutationGenerator::dfs} in \texttt{src/solver/permutations.rs}. Each minigrid is solved locally. The search respects fixed row, column, and box conflicts from the current board and masks, but it does not yet enforce compatibility with other minigrids beyond the fixed cells already present in the heuristic board.

\begin{lstlisting}[caption={Generate all minigrid permutations}]
procedure GenerateAllMinigridPermutations(board B, masks M)
    for each minigrid id i in 0 .. N - 1 do in parallel
        minigrid <- ExtractMinigrid(B, i)
        generated[i] <- GenerateMinigridPermutations(minigrid, M)
    end for

    return generated
end procedure

procedure GenerateMinigridPermutations(minigrid g, masks M)
    used <- M.boxes[g.id]
    nodes <- empty list
    payloads <- empty list

    DFSMinigrid(g, M, used, nodes, payloads)

    return GeneratedMinigrid(nodes, payloads)
end procedure
\end{lstlisting}

\begin{lstlisting}[caption={Depth-first minigrid completion with MRV}]
procedure DFSMinigrid(minigrid g, masks M, mask used, nodes, payloads)
    choice <- FindBestCell(g, M, used)

    if choice = Some(index, conflict) then
        candidates <- digits not present in conflict

        for each digit d in candidates do
            g.cells[index] <- d
            remove index from g.empty

            next_used <- used union {d}
            DFSMinigrid(g, M, next_used, nodes, payloads)

            g.cells[index] <- 0
            add index to g.empty
        end for

    else if used contains all digits 1 .. N then
        payload_id <- length(payloads)
        append copy(g.cells) to payloads
        append PermutationNode(g.cells, payload_id) to nodes

    else
        return // dead end
    end if
end procedure
\end{lstlisting}

\begin{lstlisting}[caption={Select the most constrained empty cell inside a minigrid}]
function FindBestCell(minigrid g, masks M, mask used)
    best <- None
    best_forbidden_count <- 0

    for each local index i in g.empty do
        (r, c) <- GlobalCoordinates(g.id, i)
        conflict <- M.conflict[r][c] union used

        if conflict contains all digits 1 .. N then
            return None
        end if

        forbidden_count <- CountDigits(conflict)
        if forbidden_count > best_forbidden_count then
            best <- Some(i, conflict)
            best_forbidden_count <- forbidden_count

            if forbidden_count = N - 1 then
                break
            end if
        end if
    end for

    return best
end function
\end{lstlisting}

\section{Compatibility Graph Construction}

This algorithm corresponds to \texttt{Graph::new}, \texttt{Graph::create\_edges}, and \texttt{PermutationNode} row and column compatibility in \texttt{src/types/graph}. The graph connects compatible permutations from related minigrids. Unrelated minigrids do not receive adjacency tables.

\begin{lstlisting}[caption={Build the graph container and signature indexes}]
procedure NewGraph(generated minigrids P)
    for each minigrid id i do
        graph.minigrids[i] <- P_i.nodes
        graph.payloads[i] <- P_i.payloads
        graph.row_index[i] <- BuildAxisIndex(P_i.nodes, use_rows = true)
        graph.col_index[i] <- BuildAxisIndex(P_i.nodes, use_rows = false)
    end for

    graph.pair_tables <- empty list
    graph.pair_lookup <- empty matrix
    return graph
end procedure

procedure BuildAxisIndex(nodes, use_rows)
    signatures <- empty list
    signature_to_id <- empty map
    perm_to_signature <- empty list
    signature_to_perms <- empty bitset list

    for each permutation id p and node n in nodes do
        if use_rows then
            signature <- row masks of n
        else
            signature <- column masks of n
        end if

        if signature is not in signature_to_id then
            assign next signature id
            append empty bitset to signature_to_perms
        end if

        s <- signature_to_id[signature]
        perm_to_signature[p] <- s
        add p to signature_to_perms[s]
    end for

    return AxisIndex(signatures, perm_to_signature, signature_to_perms)
end procedure
\end{lstlisting}

\begin{lstlisting}[caption={Create compatibility edges between related minigrids}]
procedure BuildCompatibilityEdges(graph G)
    clear G.pair_tables

    for left in 0 .. N - 1 do
        for right in left + 1 .. N - 1 do
            relation <- RelationshipBetween(left, right)
            if relation = NotRelated then
                continue
            end if

            if relation = RowRelated then
                L <- G.row_index[left]
                R <- G.row_index[right]
            else
                L <- G.col_index[left]
                R <- G.col_index[right]
            end if

            left_to_right <- empty table indexed by signatures of L
            right_to_left <- empty table indexed by signatures of R

            for each signature id a in L do
                for each signature id b in R do
                    if SignaturesCompatible(L.signatures[a], R.signatures[b]) then
                        left_to_right[a] <- left_to_right[a] union R.signature_to_perms[b]
                        right_to_left[b] <- right_to_left[b] union L.signature_to_perms[a]
                    end if
                end for
            end for

            append PairAdj(left, right, relation, left_to_right, right_to_left) to G.pair_tables
        end for
    end for

    G.pair_lookup <- BuildPairLookup(G.pair_tables)
end procedure
\end{lstlisting}

\begin{lstlisting}[caption={Relationship and signature compatibility}]
function RelationshipBetween(a, b)
    same_block_row <- floor(a / K) = floor(b / K)
    same_block_col <- (a mod K) = (b mod K)

    if same_block_row and same_block_col then
        return NotRelated
    else if same_block_row then
        return RowRelated
    else if same_block_col then
        return ColRelated
    else
        return NotRelated
    end if
end function

function SignaturesCompatible(signature A, signature B)
    for local_axis in 0 .. K - 1 do
        if A[local_axis] intersects B[local_axis] then
            return false
        end if
    end for

    return true
end function
\end{lstlisting}

\section{Local Support Pruning}

This algorithm corresponds to \texttt{Pruner::run\_local} in \texttt{src/solver/pruning.rs}. It removes permutations that have no currently live compatible neighbor in at least one related minigrid. This is a local consistency pass.

\begin{lstlisting}[caption={Iterative local support pruning}]
procedure LocalSupportPrune(graph G)
    for each minigrid i do
        alive[i] <- bitset containing every permutation id in G[i]
    end for

    repeat
        changed <- false

        for each minigrid i do
            for each permutation p in G[i] do
                if p is not alive[i] then
                    continue
                end if

                supported <- true

                for each minigrid j do
                    if RelationshipBetween(i, j) = NotRelated then
                        continue
                    end if

                    compatible <- CompatibleSet(G, i, p, j)
                    if compatible contains no permutation q with alive[j][q] = true then
                        supported <- false
                        break
                    end if
                end for

                if supported = false then
                    remove p from alive[i]
                    changed <- true
                end if
            end for
        end for
    until changed = false

    keep <- alive permutation ids for each minigrid
    RetainPermutations(G, keep)
    return number of removed permutations
end procedure
\end{lstlisting}

\section{Exact Support Pruning}

This algorithm corresponds to \texttt{Pruner::run} in \texttt{src/solver/pruning.rs}. It first applies local pruning, then runs the exact configuration search. A permutation survives only if some full valid configuration uses it.

\begin{lstlisting}[caption={Exact global support pruning}]
procedure ExactSupportPrune(graph G)
    pre_local_counts <- permutation counts of G

    LocalSupportPrune(G)

    current_counts <- permutation counts of compacted G
    configurations <- SearchAllConfigurations(G)

    for each minigrid i do
        supported[i] <- array of false values sized by current_counts[i]
    end for

    for each configuration C in configurations do
        for each minigrid i do
            p <- C[i] // p is an id in the post-local-prune graph
            supported[i][p] <- true
        end for
    end for

    for each minigrid i do
        keep[i] <- all current permutation ids p where supported[i][p] = true
        remap[i] <- map current kept ids to new compact ids
    end for

    removed <- total pre_local_counts - total kept count
    RetainPermutations(G, keep)

    remapped_configurations <- empty list
    for each configuration C in configurations do
        append RemapConfiguration(C, remap) to remapped_configurations
    end for

    return (removed, remapped_configurations)
end procedure
\end{lstlisting}

\begin{lstlisting}[caption={Retain permutations and rebuild graph indexes}]
procedure RetainPermutations(graph G, keep)
    for each minigrid i do
        new_nodes[i] <- empty list
        new_payloads[i] <- empty list

        for each old permutation id p in keep[i] do
            node <- G.minigrids[i][p]
            payload <- G.payloads[i][node.payload_id]
            node.payload_id <- length(new_payloads[i])
            append payload to new_payloads[i]
            append node to new_nodes[i]
        end for
    end for

    G.minigrids <- new_nodes
    G.payloads <- new_payloads
    rebuild row and column signature indexes
    BuildCompatibilityEdges(G)
end procedure
\end{lstlisting}

\section{Configuration Search}

This algorithm corresponds to \texttt{Search} in \texttt{src/solver/extraction.rs}. It is used both by exact pruning and by extraction when configurations are not already supplied.

\begin{lstlisting}[caption={Search all compatible minigrid configurations}]
procedure SearchAllConfigurations(graph G, mode)
    for each minigrid i do
        if G has zero permutations for i then
            return empty list
        end if

        domains[i] <- all permutation ids for minigrid i
        assignments[i] <- None
    end for

    solutions <- empty list
    trail <- empty stack
    Backtrack(G, domains, assignments, trail, solutions, mode)
    return solutions
end procedure
\end{lstlisting}

\begin{lstlisting}[caption={Backtracking with MRV and propagation}]
procedure Backtrack(graph G, domains, assignments, trail, solutions, mode)
    if mode has solution cap and length(solutions) >= cap then
        return
    end if

    if every minigrid is assigned then
        append assignments as complete configuration to solutions
        return
    end if

    i <- SelectNextMinigrid(domains, assignments)
    if i = None then
        return
    end if

    candidates <- all permutation ids in domains[i]
    for each permutation p in candidates do
        marker <- size(trail)

        if AssignAndPropagate(G, domains, assignments, trail, i, p) then
            Backtrack(G, domains, assignments, trail, solutions, mode)
        end if

        UndoTo(trail, domains, assignments, marker)
    end for
end procedure

function SelectNextMinigrid(domains, assignments)
    best <- None
    best_count <- infinity

    for each minigrid i do
        if assignments[i] is set then
            continue
        end if

        count <- size(domains[i])
        if count < best_count then
            best <- i
            best_count <- count
        end if

        if best_count = 0 then
            break
        end if
    end for

    return best
end function
\end{lstlisting}

\begin{lstlisting}[caption={Assignment, propagation, and trail-based undo}]
procedure AssignAndPropagate(graph G, domains, assignments, trail, i, p)
    queue <- [(i, p)]
    SaveState(trail, domains, assignments, i)
    assignments[i] <- p
    domains[i] <- singleton set {p}

    while queue is not empty do
        (current_i, current_p) <- pop(queue)

        for each minigrid j do
            if RelationshipBetween(current_i, j) = NotRelated then
                continue
            end if

            compatible <- CompatibleSet(G, current_i, current_p, j)

            if assignments[j] = Some(q) then
                if q is not in compatible then
                    return false
                end if
                continue
            end if

            SaveState(trail, domains, assignments, j)
            changed <- Intersect(domains[j], compatible)

            if domains[j] is empty then
                return false
            end if

            if changed and size(domains[j]) = 1 then
                q <- only element of domains[j]
                assignments[j] <- q
                push (j, q) into queue
            end if
        end for
    end while

    return true
end procedure

procedure UndoTo(trail, domains, assignments, marker)
    while size(trail) > marker do
        entry <- pop(trail)
        domains[entry.minigrid] <- entry.domain
        assignments[entry.minigrid] <- entry.assignment
    end while
end procedure
\end{lstlisting}

\section{Solution Extraction and Classification}

This algorithm corresponds to \texttt{Extractor::run\_with\_configurations}, \texttt{Extractor::reconstruct}, and the classification logic in \texttt{src/solver/mod.rs}.

\begin{lstlisting}[caption={Extract and validate solved boards}]
procedure ExtractSolutions(graph G, configurations)
    solutions <- empty list

    for each configuration C in configurations do
        board <- ReconstructBoard(G, C)

        if IsValidSudokuBoard(board) then
            append Solution(board, C) to solutions
        end if
    end for

    return solutions
end procedure

procedure ReconstructBoard(graph G, configuration C)
    initialize cells[N][N] with zeros

    for each minigrid id i do
        p <- C[i]
        perm_cells <- G.permutation_cells(i, p)
        base_row <- floor(i / K) * K
        base_col <- (i mod K) * K

        for local index t in 0 .. N - 1 do
            local_row <- floor(t / K)
            local_col <- t mod K
            row <- base_row + local_row
            col <- base_col + local_col
            cells[row][col] <- perm_cells[t]
        end for
    end for

    return Board(cells)
end procedure

function Classify(solutions)
    if length(solutions) = 0 then
        return Unsolvable
    else if length(solutions) = 1 then
        return Unique
    else
        return Ambiguous(length(solutions))
    end if
end function
\end{lstlisting}

\section{Algorithm Map to Source Files}

\begin{itemize}
    \item Complete pipeline: \texttt{src/solver/mod.rs}
    \item Mask initialization: \texttt{src/types/masks.rs}
    \item Hidden singles: \texttt{src/solver/heuristics.rs}
    \item Minigrid permutation generation: \texttt{src/solver/permutations.rs}
    \item Graph construction and compatibility: \texttt{src/types/graph/mod.rs}, \texttt{src/types/graph/node.rs}
    \item Local and exact pruning: \texttt{src/solver/pruning.rs}
    \item Configuration search and reconstruction: \texttt{src/solver/extraction.rs}
\end{itemize}

\end{document}
